% !TeX root = ../教材题库_全部条目_转写工作稿.tex
% Source: XJ-B1-C03-HW-P078-Y1137
% 第3章 / 3.1 函数 / 教材 p.78 / PDF 第85页
% OCR未校对：公式与文字均可能有误；未完成数学TeX转写；禁止直接发给学生
% 原题及相邻原解截图（核对用）：../题目原图/XJ-B1-C03-HW-P078-Y1137-01.webp
% 本文件只在转写完成后作为题干使用；OCR 条目尚不是可靠题干。
\begin{bankitem}{习题3.1 第15题}{第3章 / 3.1 函数 / 教材 p.78 / PDF 第85页}{XJ-B1-C03-HW-P078-Y1137}{OCR未校对：公式与文字均可能有误；未完成数学TeX转写；禁止直接发给学生}
\begin{ocrdraft}
[教材p.78 / PDF第85页]
\par
巧．如图，四边形ABCD是边长为4的正方形，一个动点P从点A出发，沿
\par
着边AB—BC—CD—DA运动，返回到点A后停止运动，设点P走过的路程为
\par
（1）将线段AP的长度表示成的函数丆l(x）；
\par
（2）设点0是正方形ABCD的中心，当点A，0，P三点不共线时，将AAOP
\par
的面积表示成的函数丆20）．
\end{ocrdraft}
\end{bankitem}
